% =====================================================================
% aestiva.tex -- LaTeX conversion of aestiva.md (Breviarium Romanum,
% Pars Aestiva), BODY CONTENT ONLY.
%
% No preamble, \documentclass, or \newcommand definitions are included
% here, per instruction -- the command vocabulary below follows the
% same legend established for hiemalis.tex and verna.tex (one-column
% vs. two-column, red vs. black, etc.).
%
% BATCH 1 of N: Ordinarium Divini Officii, Ad Matutinum and Ad Laudes
% (aestiva.md lines 1-463)
% =====================================================================

% [ornamental image beginning Pars Aestiva]
\partheading{Ordinarium Divini Officii}
\partsubtitle{Juxta Ritum Romanum Persolvendi}

\rubric{Clericis in sacris constitútis, qui íntegrum divínum Offícium, quamvis in partes distribútum, coram Ssmo Sacraménto, sive públicæ adoratióni expósito sive in tabernáculo adserváto, devóte recitáverint, concéditur: Indulgéntia plenária, si prætérea sacramentálem confessiónem instítuerint, ad eucharísticum Convívium accésserint et ad mentem Summi Pontíficis oráverint. \textit{Pius Pp. XI, 23 Octóbris 1930.}}

\rubric{Iis vero, qui divínum Offícium coram Ssmo Sacraménto, ut supra, ex parte tantum recitáverint, concéditur: Indulgéntia quingentórum diérum pro unaquáque hora canónica. \textit{Pius Pp. XI, 18 Maji 1933.}}

\rubric{Iisdem autem cléricis, qui divínum Offícium in álias preces commutátum rite obtinúerint, si preces ejúsmodi coram Ssmo Sacraménto, uti supra, devóte recitáverint, concéditur: Indulgéntia plenária, addítis sacramentáli confessióne, sacra Communióne et oratióne ad mentem Summi Pontíficis. \textit{Pius Pp. XI, 7 Novémbris 1932.}}

\sectionheading{Ante Divinum Officium}

\rubric{Antequam inchoëtur Offícium laudabíliter dícitur, sub singulári semper número, séquens Orátio; pro qua Summus Póntifex Pius Pp. XI sub die 17 Novémbris 1933 Indulgéntiam trium annórum concéssit.}

\oratioheading{Oratio}

\lettrine[lines=2]{A}{PERI}, Dómine, os meum ad benedicéndum nomen sanctum tuum: munda quoque cor meum ab ómnibus vanis, pervérsis et aliénis cogitatiónibus; intelléctum illúmina, afféctum inflámma, ut digne, atténte ac devóte hoc Offícium recitáre váleam, et exaudíri mérear ante conspéctum divínæ Majestátis tuæ. Per Christum, Dóminum nostrum. \R\ Amen.

Dómine, in unióne illíus divínæ intentiónis, qua ipse in terris laudes Deo persolvísti, has tibi Horas (\textit{vel} hanc tibi Horam) persólvo.

\sectionheading{Ante singulas Horas}

\rubric{Ante Matutínum, sicut et ante Laudes a Matutíno extra Chorum separátas, et ante omnes Horas, præter quam ad Completórium, dícitur secréto:}

\lettrine[lines=2]{P}{ATER} noster, qui es in cælis, sanctificétur nomen tuum. Advéniat regnum tuum. Fiat volúntas tua, sicut in cælo et in terra. Panem nostrum quotidiánum da nobis hódie. Et dimítte nobis débita nostra, sicut et nos dimíttimus debitóribus nostris. Et ne nos indúcas in tentatiónem: sed líbera nos a malo. Amen.

\lettrine[lines=2]{A}{VE}, María, grátia plena; Dóminus tecum: benedícta tu in muliéribus, et benedíctus fructus ventris tui Jesus. Sancta María, Mater Dei, ora pro nobis peccatóribus, nunc et in hora mortis nostræ. Amen.

\rubric{In princípio autem Matutíni ac Primæ additur secréto Symbolum Apostolórum:}

\lettrine[lines=2]{C}{REDO} in Deum, Patrem omnipoténtem, Creatórem cæli et terræ. Et in Jesum Christum Fílium ejus únicum, Dóminum nostrum: qui concéptus est de Spíritu Sancto, natus ex María Vírgine, passus sub Póntio Piláto, crucifíxus, mórtuus, et sepúltus: descéndit ad ínferos; tértia die resurréxit a mórtuis; ascéndit ad cælos; sedet ad déxteram Dei Patris omnipoténtis: inde ventúrus est judicáre vivos et mórtuos. Credo in Spíritum Sanctum, sanctam Ecclésiam cathólicam, Sanctórum communiónem, remissiónem peccatórum, carnis resurrectiónem, vitam ætérnam. Amen.

\rubric{Officium ad Matutinum et ad reliquas Horas, nisi aliter suo loco indicetur, per totum annum persolvitur juxta formam, quæ in Rubricis hujus Ordinarii exhibetur.}

\hourheading{Ad Matutinum}

Pater noster, Ave María et Credo.

\rubric{Deinde, clara voce, dicitur Versus:}

\V\ \lettrine[lines=2]{D}{OMINE}, lábia mea apéries. \R\ Et os meum annuntiábit laudem tuam. \V\ Deus, in adjutórium meum inténde. \R\ Dómine, ad adjuvándum me festína.

Glória Patri, et Fílio, et Spirítui Sancto. Sicut erat in princípio, et nunc, et semper, et in sæcula sæculórum. Amen. Allelúja.

\rubric{Postea dicitur conveniens Invitatorium, quod ante Psalmum bis integrum recitatur, et ad singulos ejusdem Psalmi versus vel integrum vel dimidiatum ab Asterisco * alternis vicibus, ut infra, repetitur.}

\psalmheading{Psalmus 94}

\lettrine[lines=2]{V}{ENITE}, exsultémus Dómino, jubilémus Deo, salutári nostro: præoccupémus fáciem ejus in confessióne, et in psalmis jubilémus ei.

\rubric{Repetitur integrum Invitatorium.}

Quóniam Deus magnus Dóminus, et Rex magnus super omnes deos: quóniam non repéllet Dóminus plebem suam: quia in manu ejus sunt omnes fines terræ, et altitúdines móntium ipse cónspicit.

\rubric{Repetitur altera pars Invitatorii.}

\rubric{In sequenti Psalmi versu, ad verba veníte, adorémus, et procidámus ante Deum, genuflectitur.}

Quóniam ipsíus est mare, et ipse fecit illud, et áridam fundavérunt manus ejus: veníte, adorémus, et procidámus ante Deum: plorémus coram Dómino, qui fecit nos, quia ipse est Dóminus, Deus noster; nos autem pópulus ejus, et oves páscuæ ejus.

\rubric{Repetitur integrum Invitatorium.}

Hódie, si vocem ejus audíeritis, nolíte obduráre corda vestra, sicut in exacerbatióne secúndum diem tentatiónis in desérto: ubi tentavérunt me patres vestri, probavérunt et vidérunt ópera mea.

\rubric{Repetitur altera pars Invitatorii.}

Quadragínta annis próximus fui generatióni huic, et dixi: Semper hi errant corde; ipsi vero non cognovérunt vias meas: quibus jurávi in ira mea: Si introíbunt in réquiem meam.

\rubric{Repetitur integrum Invitatorium.}

Glória Patri, et Fílio, et Spirítui Sancto. Sicut erat in princípio, et nunc, et semper, et in sæcula sæculórum. Amen.

\rubric{Repetitur altera pars Invitatorii.}

\rubric{Denique repetitur integrum Invitatorium.}

\rubric{Expleto Psalmo, dicitur Hymnus Invitatorio respondens.}

\variantheading{In Officio dominicali}

\rubric{In omnibus Dominicis, extra Octavas Ssmi Corporis Christi ac Sacrat. Cordis Jesu, Invitatorium et Hymnus ut in Psalterio.}

\variantheading{In Officio feriali}

\rubric{In omnibus Feriis, Invitatorium et Hymnus de Feria currenti, ut in Psalterio.}

\variantheading{In Festis}

\rubric{In Officio cujuslibet Festi vel Octavæ, ac S. Mariæ in Sabbato, Invitatorium et Hymnus ut in Proprio vel Communi; in Dominicis autem infra Octavas Ssmi Corporis Christi ac Sacratissimi Cordis Jesu, si de eis fiat Officium, Invitatorium et Hymnus ut in Proprio de Tempore. Quæ norma ad reliquas item Horas erit observanda pro illis Officii partibus, quæ juxta Rubricas ex Ordinario vel Psalterio non fuerint in Festis desumendæ.}

\variantheading{In Officio novem Lectionum}

\rubric{Expleto Hymno, dicuntur Antiphonæ convenientes, quæ in Officiis ritus duplicis ante et post Psalmos integræ recitantur; in Officiis autem ritus semiduplicis initio Psalmi inchoantur tantum et usque ad Asteriscum * perducuntur, atque in fine integræ pronuntiantur; servata alterutra norma etiam in omnibus Antiphonis, quæ ad Laudes ac Vesperas adhibentur.}

\nocturnoheading{In I Nocturno}

\rubric{Sub congruentibus Antiphonis dicuntur tres Psalmi, ac deinde subjungitur Versus, prouti Officium occurrens requirit.}

\rubric{Post Versum cujuslibet Nocturni dicitur}

Pater noster \rubric{secreto usque ad} \V\ Et ne nos indúcas in tentatiónem. \R\ Sed líbera nos a malo.

\absolutioheading{Absolutio}

\lettrine[lines=2]{E}{XÁUDI}, Dómine Jesu Christe, preces servórum tuórum, et miserére nobis: Qui cum Patre et Spíritu Sancto vivis et regnas in sæcula sæculórum. \R\ Amen.

\V\ Jube, domne, benedícere.

\rubric{¶ Extra Chorum, quando ab uno tantum recitatur Officium, ante singulas Lectiones Matutini atque ad Lectionem brevem Primæ et Completorii, dicitur: Jube, Dómine, benedícere; et subjungitur congruens Benedictio. Ab Episcopo autem, ultimam Matutini Lectionem cantaturo, item dicitur: Jube, Dómine, benedícere; et respondetur a Choro: Amen.}

\variantheading{Pro I Lectione}

\benedlabel{Bened.} Benedictióne perpétua benedícat nos Pater ætérnus. \R\ Amen.

\rubric{Deinde dicuntur in unoquoque Nocturno Lectiones, prouti Officium occurrens requirit, et in fine cujuslibet Lectionis additur:}

\V\ Tu autem, Dómine, miserére nobis. \R\ Deo grátias.

\rubric{Post quamlibet vero Lectionem, quæ Hymnum Te Deum immediate non præcedat, congruens dicitur Responsorium, et in fine ultimi Responsorii cujusque Nocturni additur Versus: Glória Patri, et Fílio, et Spirítui Sancto, et Responsorium ab Asterisco *, et quidem a secundo, si duo fuerint, repetitur.}

\variantheading{Pro II Lectione}

\benedlabel{Bened.} Unigénitus Dei Fílius nos benedícere et adjuváre dignétur. \R\ Amen.

\variantheading{Pro III Lectione}

\benedlabel{Bened.} Spíritus Sancti grátia illúminet sensus et corda nostra. \R\ Amen.

\nocturnoheading{In II Nocturno}

\rubric{Sub congruentibus item Antiphonis dicuntur tres Psalmi et Versus, sicut in I Nocturno.}

\rubric{Post Versum dicitur.}

Pater noster \rubric{secreto usque ad} \V\ Et ne nos indúcas in tentatiónem. \R\ Sed líbera nos a malo.

\absolutioheading{Absolutio}

\lettrine[lines=2]{I}{PSIUS} píetas et misericórdia nos ádjuvet, qui cum Patre et Spíritu Sancto vivit et regnat in sæcula sæculórum. \R\ Amen.

\V\ Jube, domne, benedícere.

\variantheading{Pro IV Lectione}

\benedlabel{Bened.} Deus Pater omnípotens sit nobis propítius et clemens. \R\ Amen.

\variantheading{Pro V Lectione}

\benedlabel{Bened.} Christus perpétuæ det nobis gáudia vitæ. \R\ Amen.

\variantheading{Pro VI Lectione}

\benedlabel{Bened.} Ignem sui amóris accéndat Deus in córdibus nostris. \R\ Amen.

\nocturnoheading{In III Nocturno}

\rubric{Sub congruentibus denique Antiphonis dicuntur tres Psalmi et Versus, sicut in I et II Nocturno.}

\rubric{Post Versum dicitur}

Pater noster \rubric{secreto usque ad} \V\ Et ne nos indúcas in tentatiónem. \R\ Sed líbera nos a malo.

\absolutioheading{Absolutio}

\lettrine[lines=2]{A}{} \textsc{vinculis} peccatórum nostrórum absólvat nos omnípotens et miséricors Dóminus. \R\ Amen.

\V\ Jube, domne, benedícere.

\variantheading{Pro VII Lectione}

\benedlabel{Bened.} Evangélica léctio sit nobis salus et protéctio. \R\ Amen.

\variantheading{Pro VIII Lectione}

\rubric{In Festis Domini et in Dominicis:}

\benedlabel{Bened.} Divínum auxílium máneat semper nobíscum. \R\ Amen.

\rubric{In Festis beatæ Mariæ Virginis:}

\benedlabel{Bened.} Cujus festum cólimus, ipsa Virgo virgínum intercédat pro nobis ad Dóminum. \R\ Amen.

\rubric{In Festis Sanctorum:}

\benedlabel{Bened.} Cujus (\textit{vel} Quorum \textit{aut} Quarum) festum cólimus, ipse (\textit{vel} ipsa \textit{aut} ipsi \textit{vel} ipsæ) intercédat (\textit{vel} intercédant) pro nobis ad Dóminum. \R\ Amen.

\variantheading{Pro IX Lectione}

\benedlabel{Bened.} Ad societátem cívium supernórum perdúcat nos Rex Angelórum. \R\ Amen.

\rubric{Si autem legenda sit Lectio ix de Homilia cum Evangelio Dominicæ, vel Feriæ, aut Vigiliæ:}

\benedlabel{Bened.} Per evangélica dicta deleántur nostra delícta. \R\ Amen.

\variantheading{In Officio trium Lectionum}

\nocturnoheading{In Nocturno}

\rubric{In Officiis trium Lectionum, post Hymnum dicuntur Antiphonæ convenientes, quæ initio Psalmi inchoantur tantum et usque ad Asteriscum * perducuntur, ac deinde in fine integræ pronuntiantur.}

\rubric{Sub eisdem vero Antiphonis dicuntur novem Psalmi Feriæ currentis, quibus subjungitur Versus in III Nocturno positus, omissis Versibus pro I et II Nocturno assignatis.}

\rubric{Post Versum dicitur}

Pater noster \rubric{secreto usque ad} \V\ Et ne nos indúcas in tentatiónem. \R\ Sed líbera nos a malo.

\absolutioheading{Absolutio}

\rubric{Feria II et V}

Exáudi, Dómine, \rubric{ut in I Nocturno.}

\rubric{Feria III et VI}

Ipsíus píetas, \rubric{ut in II Nocturno.}

\rubric{Feria IV et Sabbato}

A vínculis, \rubric{ut in III Nocturno Officii novem Lectionum.}

\rubric{Deinde leguntur Lectiones cum Responsoriis, prouti Officium occurrens requirit; et ante eas dicuntur sequentes}

\variantheading{Benedictiones}

\rubric{In Feriis, quando legitur Homilia cum Evangelio, erit:}

\benedlabel{1 Bened.} Evangélica léctio sit nobis salus et protéctio. \R\ Amen.

\benedlabel{2 Bened.} Divínum auxílium máneat semper nobíscum. \R\ Amen.

\benedlabel{3 Bened.} Ad societátem cívium supernórum perdúcat nos Rex Angelórum. \R\ Amen.

\rubric{In Feriis, quando non legitur Homilia cum Evangelio, Feria II et V Benedictiones ut in I Nocturno, Feria III et VI ut in II Nocturno Officii novem Lectionum; Feria autem IV et Sabbato erit:}

\benedlabel{1 Bened.} Ille nos benedícat, qui sine fine vivit et regnat. \R\ Amen.

\benedlabel{2 Bened.} Divínum auxílium máneat semper nobíscum. \R\ Amen.

\benedlabel{3 Bened.} Ad societátem cívium supernórum perdúcat nos Rex Angelórum. \R\ Amen.

\rubric{In Festis Sanctorum erit:}

\benedlabel{1 Bened.} Ille nos benedícat, qui sine fine vivit et regnat. \R\ Amen.

\benedlabel{2 Bened.} Cujus (\textit{vel} Quorum \textit{aut} Quarum) festum cólimus, ipse (\textit{vel} ipsa \textit{aut} ipsi \textit{vel} ipsæ) intercédat (\textit{vel} intercédant) pro nobis ad Dóminum. \R\ Amen.

\benedlabel{3 Bened.} Ad societátem cívium supernórum perdúcat nos Rex Angelórum. \R\ Amen.

\rubric{In Officio de sancta María in Sábbato Absolútio et Benedictiónes, quæ ibídem ponúntur.}

\rubric{Post últimam Lectiónem, in ómnibus Domínicis per Annum minóribus, in Festis cujúsvis ritus, per omnes Octávas, et in Offício sanctæ Maríæ in Sábbato, dícitur Hymnus Ambrosiánus. Et numquam dícitur in feriáli Offício. Quando vero Hymnus prædíctus omíttitur, ejus loco dícitur IX aut III Responsórium.}

\hymnusheading{Hymnus Ambrosianus}

\lettrine[lines=2]{T}{E} Deum laudámus: * te Dóminum confitémur.

Te ætérnum Patrem * omnis terra venerátur.

Tibi omnes Angeli, * tibi Cæli, et univérsæ Potestátes:

Tibi Chérubim et Séraphim * incessábili voce proclámant:

Sanctus, * Sanctus, * Sanctus * Dóminus, Deus Sábaoth.

Pleni sunt cæli et terra * majestátis glóriæ tuæ.

Te gloriósus * Apostolórum chorus,

Te Prophetárum * laudábilis númerus,

Te Mártyrum candidátus * laudat exércitus.

Te per orbem terrárum * sancta confitétur Ecclésia,

Patrem * imménsæ majestátis;

Venerándum tuum verum * et únicum Fílium;

Sanctum quoque * Paráclitum Spíritum.

Tu Rex glóriæ, * Christe.

Tu Patris * sempitérnus es Fílius.

Tu, ad liberándum susceptúrus hóminem, * non horruísti Vírginis úterum.

Tu, devícto mortis acúleo, * aperuísti credéntibus regna cælórum.

Tu ad déxteram Dei sedes, * in glória Patris.

Judex créderis * esse ventúrus.

\rubric{¶ Sequens versus dicitur flexis genibus.}

Te ergo quæsumus, tuis fámulis súbveni, * quos pretióso sánguine redemísti.

Ætérna fac cum Sanctis tuis * in glória numerári.

Salvum fac pópulum tuum, Dómine, * et bénedic hereditáti tuæ.

Et rege eos, * et extólle illos usque in ætérnum.

Per síngulos dies * benedícimus te;

Et laudámus nomen tuum in séculum, * et in séculum séculi.

Dignáre, Dómine, die isto * sine peccáto nos custodíre.

Miserére nostri, Dómine, * miserére nostri.

Fiat misericórdia tua, Dómine, super nos, * quemádmodum sperávimus in te.

In te, Dómine, sperávi: * non confúndar in ætérnum.

\rubric{Dicto Te Deum, aut ultimo Responsorio, statim incipiuntur Laudes a Versu Deus, in adjutórium, ut infra.}

\rubric{In privata vero recitatione Matutinum a Laudibus separari potest; quo in casu, post Hymnum Te Deum, vel post ultimum Responsorium, dicitur:}

\V\ Dóminus vobíscum. \R\ Et cum spíritu tuo.

\rubric{¶ Hic Versus Dóminus vobíscum non dicitur ab eo, qui non est saltem in ordine Diaconatus; sed ejus loco substituitur, etiam in aliis Horis, quoties in antecedentibus Precibus dictus non fuerit:}

\V\ Dómine, exáudi oratiónem meam. \R\ Et clamor meus ad te véniat.

Orémus.

\rubric{Et dicitur Oratio ut ad Laudes.}

\rubric{Post Orationem additur:}

\V\ Dóminus vobíscum. \R\ Et cum spíritu tuo.

\V\ Benedicámus Dómino. \R\ Deo grátias.

\V\ Fidélium ánimæ per misericórdiam Dei requiéscant in pace. \R\ Amen.

\rubric{Deinde dicitur Pater noster totum secreto.}

\rubric{In publica vero recitatione Matutinum in Choro a Laudibus separari non potest.}

\hourheading{Ad Laudes}

\rubric{Si Laudes extra Chorum separentur a Matutino, ante eas dicitur secreto Pater noster et Ave María; secus absolute incipiuntur, ut sequitur:}

\V\ Deus, in adjutórium.

\rubric{Deinde, sub congruentibus Antiphonis, dicuntur Psalmi et Canticum, prouti Officium occurrens requirit.}

\rubric{Repetita vero post ultimum Psalmum Antiphona, dicuntur Capitulum, Hymnus et Versus, prouti Officium occurrens exigit.}

\variantheading{In Officio dominicali}

\rubric{In omnibus Dominicis, extra Octavas Ssmi Corporis Christi ac Sacratissimi Cordis Jesu, Capitulum, Hymnus et Versus ut in Psalterio.}

\variantheading{In Officio feriali}

\rubric{In omnibus Feriis, Capitulum, Hymnus et Versus de Feria currenti, ut in Psalterio.}

\rubric{Recitato Versu, dicitur cum Antiphona convenienti sequens}

\canticumheading{Canticum Zachariæ}{Luc. 1, 68-79}

\lettrine[lines=2]{B}{ENEDICTUS} \textdagger\ Dóminus, Deus Israël, * quia visitávit, et fecit redemptiónem plebis suæ:

Et eréxit cornu salútis nobis: * in domo David, púeri sui.

Sicut locútus est per os sanctórum, * qui a sæculo sunt, Prophetárum ejus:

Salútem ex inimícis nostris, * et de manu ómnium, qui odérunt nos:

Ad faciéndam misericórdiam cum pátribus nostris: * et memorári testaménti sui sancti.

Jusjurándum, quod jurávit ad Abraham, patrem nostrum, * datúrum se nobis:

Ut sine timóre, de manu inimicórum nostrórum liberáti, * serviámus illi.

In sanctitáte, et justítia coram ipso, * ómnibus diébus nostris.

Et tu, puer, Prophéta Altíssimi vocáberis: * præíbis enim ante fáciem Dómini paráre vias ejus:

Ad dandam sciéntiam salútis plebi ejus: * in remissiónem peccatórum eórum:

Per víscera misericórdiæ Dei nostri: * in quibus visitávit nos, óriens ex alto:

Illumináre his, qui in ténebris, et in umbra mortis sedent: * ad dirigéndos pedes nostros in viam pacis.

Glória Patri, et Fílio, * et Spirítui Sancto.

Sicut erat in princípio, et nunc, et semper, * et in sæcula sæculórum. Amen.

\rubric{¶ Hic Versus Glória Patri dicitur in fine omnium Canticorum, nisi aliter adnotetur.}

\rubric{Expleto Cantico ac repetita Antiphona, si Preces non fuerint recitandæ, statim dicitur Versus Dóminus vobíscum cum Oratione et reliquis omnibus usque ad finem Horæ, ut infra.}

\rubric{In omnibus vero Vigiliis, si fiat Officium de Vigilia currenti, post repetitam ad Benedictus Antiphonam dicuntur flexis genibus sequentes Preces, quæ aliis temporibus omittuntur:}

KYRIE, eléison. Christe, eléison. Kýrie, eléison.

\rubric{Pater noster, quod dicitur a solo Hebdomadario totum clara voce, usque ad} \V\ Et ne nos indúcas in tentatiónem. \R\ Sed líbera nos a malo.

\V\ Ego dixi: Dómine, miserére mei. \R\ Sana ánimam meam, quia peccávi tibi. \V\ Convértere, Dómine, úsquequo? \R\ Et deprecábilis esto super servos tuos. \V\ Fiat misericórdia tua, Dómine, super nos. \R\ Quemádmodum sperávimus in te. \V\ Sacerdótes tui induántur justítiam. \R\ Et sancti tui exsúltent.

\V\ Orémus pro beatíssimo Papa nostro N. \R\ Dóminus consérvet eum, et vivíficet eum, et beátum fáciat eum in terra, et non tradat eum in ánimam inimicórum ejus.

\V\ Orémus et pro Antístite nostro N. \R\ Stet et pascat in fortitúdine tua, Dómine, in sublimitáte nóminis tui.

\rubric{¶ Romæ præcedens Versus cum suo Responsorio omittitur; alibi vero, ad litteram N., ab omnibus nomen diœcesani Episcopi exprimatur. Vacante Apostolica vel Episcopali Sede, alteruter vel uterque respectivus Versus cum suo Responsorio prætéritur.}

\V\ Dómine, salvum fac regem. \R\ Et exáudi nos in die, qua invocavérimus te. \V\ Salvum fac pópulum tuum, Dómine, et bénedic hereditáti tuæ. \R\ Et rege eos, et extólle illos usque in ætérnum. \V\ Meménto Congregatiónis tuæ. \R\ Quam possedísti ab inítio. \V\ Fiat pax in virtúte tua. \R\ Et abundántia in túrribus tuis. \V\ Orémus pro benefactóribus nostris. \R\ Retribúere dignáre, Dómine, ómnibus, nobis bona faciéntibus propter nomen tuum, vitam ætérnam. Amen. \V\ Orémus pro fidélibus defúnctis. \R\ Réquiem ætérnam dona eis, Dómine, et lux perpétua lúceat eis. \V\ Requiéscant in pace. \R\ Amen. \V\ Pro frátribus nostris abséntibus. \R\ Salvos fac servos tuos, Deus meus, sperántes in te. \V\ Pro afflíctis et captívis. \R\ Líbera eos, Deus Israël, ex ómnibus tribulatiónibus suis. \V\ Mitte eis, Dómine, auxílium de sancto. \R\ Et de Sion tuére eos. \V\ Dómine, Deus virtútum, convérte nos. \R\ Et osténde fáciem tuam, et salvi érimus. \V\ Exsúrge, Christe, ádjuva nos. \R\ Et líbera nos propter nomen tuum. \V\ Dómine, exáudi oratiónem meam. \R\ Et clamor meus ad te véniat.

\rubric{Absolutis Precibus, vel, si Preces locum non habeant, repetita Antiphona post Benedictus, statim dicitur:}

\V\ Dóminus vobíscum. \R\ Et cum spíritu tuo.

Orémus.

\rubric{Et dicitur Oratio conveniens. Postea fiunt Commemorationes, si occurrant.}

\rubric{Deinde in Officio semiduplici aut simplici tam de Tempore quam de Sanctis, præter omnes dies, in quibus habita sit Commemoratio Duplicis, vel factum fuerit Officium aut Commemoratio cujusvis Octavæ, fit sequens}

\suffragiumheading{Suffragium de omnibus Sanctis}

\antlabel{Ant.} Beáta Dei Génitrix Virgo María, Sanctíque omnes intercédant pro nobis ad Dóminum.

\V\ Mirificávit Dóminus Sanctos suos. \R\ Et exaudívit eos clamántes ad se.

Orémus. \oratiolabel{Oratio}

\lettrine[lines=2]{A}{} \textsc{cunctis} nos, quæsumus, Dómine, mentis et córporis defénde perículis: et, intercedénte beáta et gloriósa semper Vírgine Dei Genitríce María, cum beáto Joseph, beátis Apóstolis tuis Petro et Paulo, atque beáto N. et ómnibus Sanctis, salútem nobis tríbue benígnus et pacem; ut, destrúctis adversitátibus et erróribus univérsis, Ecclésia tua secúra tibi sérviat libertáte. Per eúmdem Dóminum nostrum. \R\ Amen.

\rubric{¶ In hac Oratione, ad litteram N., exprimitur nomen Titularis propriæ Ecclesiæ, dummodo Titulus non sit Persona divina vel Mysterium Domini, aut de ipso ratione Vigiliæ vel Festi factum non fuerit Officium vel Commemoratio, aut illius nomen non sit in eadem Suffragii Oratione expressum; ac nomina sanctorum Angelorum et sancti Joannis Baptistæ, si Titulares fuerint, præponuntur nomini sancti Joseph. In hisce autem omnibus casibus omittuntur verba atque beáto N.}

\rubric{Si vero Officium vel Commemoratio fuerit de beáta María Vírgine, tunc dicitur ut infra; atque insuper, sicubi Titularis juxta superiorem Rubricam ante omnes fuerit nominandus, ad Orationem particula cum omittitur ante nomen sanctorum Petri et Pauli ac præponitur nomini sancti Joseph.}

\antlabel{Ant.} Sancti omnes intercédant pro nobis ad Dóminum.

\V\ Mirificávit Dóminus Sanctos suos. \R\ Et exaudívit eos clamántes ad se.

Orémus. \oratiolabel{Oratio}

\lettrine[lines=2]{A}{} \textsc{cunctis} nos, quæsumus, Dómine, mentis et córporis defénde perículis: et, intercedénte beáto Joseph, cum beátis Apóstolis tuis Petro et Paulo, atque beáto N. et ómnibus Sanctis, salútem nobis tríbue benígnus et pacem; ut, destrúctis adversitátibus et erróribus univérsis, Ecclésia tua secúra tibi sérviat libertáte. Per Dóminum. \R\ Amen.

\rubric{Post ultimam Orationem additur:}

\V\ Dóminus vobíscum. \R\ Et cum spíritu tuo.

\V\ Benedicámus Dómino. \R\ Deo grátias.

\V\ Fidélium ánimæ per misericórdiam Dei requiéscant in pace. \R\ Amen.

\rubric{Deinde, in publica Officii recitatione, si discedendum sit a Choro, itemque in recitatione privata, si tunc terminetur Officium; alioquin, in fine ultimæ Horæ dicitur Pater noster, totum secreto; eoque recitato, subjungitur:}

\V\ Dóminus det nobis suam pacem. \R\ Et vitam ætérnam. Amen.

\rubric{Et immediate dicitur, cum suis Versu et Oratione, finalis beatæ Mariæ Virginis Antiphona, quæ inferius post Completorium assignatur.}

\rubric{Postea concluditur:}

\V\ Divínum auxílium máneat semper nobíscum. \R\ Amen.

% ===== END OF BATCH 1 -- Covers aestiva.md lines 1-463: the opening
% of the Ordinarium Divini Officii (its introductory indulgence
% rubrics and the "Aperi, Dómine" Oratio), "Ante singulas Horas" (the
% Pater noster, Ave María, and Credo spelled out in full for the
% first and only time -- each given a lettrine since the source
% presents them in full capitals as complete standalone prayers,
% unlike their thousands of later plain-text abbreviated recurrences
% throughout the volume), "Ad Matutinum" (the Invitatorium Psalm 94,
% the three-Nocturn apparatus with its Absolutios, Benedictiones, and
% "Pro N Lectione" headings, the Officium trium Lectionum variant,
% and the Te Deum under its customary heading "Hymnus Ambrosianus"),
% and "Ad Laudes" through the Suffragium de omnibus Sanctis (in both
% its Marian-titular and Joseph-titular forms) and the closing
% versicles. This Pars Aestiva Ordinarium is notably simpler than its
% Pars Verna counterpart (Batch 1 of verna.tex): since the Summer
% volume's Proprium de Tempore never touches Lent, Passiontide, or
% Paschaltide, this copy of the Ordinarium omits the "Tempore
% Quadragesimæ"/"Tempore Passionis" Invitatorium-Hymnus variants and
% the Paschaltide "Commemoratio de Cruce" that verna.md's Ordinarium
% carried, confirming it is independently set, not a mechanical copy
% of the other volumes' shared material -- every line here has been
% transcribed fresh from aestiva.md's own text (which also lacks the
% stray "adservátο" Greek-omicron typo found in verna.md's otherwise
% near-identical opening rubric, confirming independent proofing).
% Roman numerals in the "Pro N Lectione" headings are capitalized
% (e.g. \variantheading{Pro I Lectione}) though the source itself
% gives them lowercase ("Pro i Lectione"), matching the established
% normalization already present throughout hiemalis.tex and
% verna.tex for this specific heading type. The document continues
% with "Ad Primam" (aestiva.md line 465) =====
\hourheading{Ad Primam}

Pater, Ave et Credo.

\V\ Deus, in adjutórium.

\hymnusheading{Hymnus}

\lettrine[lines=2]{J}{AM} lucis orto sídere,\\
Deum precémur súpplices,\\
Ut in diúrnis áctibus\\
Nos servet a nocéntibus.

Linguam refrénans témperet,\\
Ne litis horror ínsonet:\\
Visum fovéndo cóntegat,\\
Ne vanitátes háuriat.

Sint pura cordis íntima,\\
Absístat et vecórdia:\\
Carnis terat supérbiam\\
Potus cibíque párcitas.

Ut, cum dies abscésserit\\
Noctémque sors redúxerit,\\
Mundi per abstinéntiam\\
Ipsi canámus glóriam.

Deo Patri sit glória,\\
Ejúsque soli Fílio,\\
Cum Spíritu Paráclito\\
Nunc et per omne sæculum. Amen.

\rubric{¶ Conclusio communis, in Hymno præcedenti, et in aliis idem metrum habentibus, semper omittitur, quando specialis in omnibus Horis adhibenda præscribitur; et, si plures Conclusiones propriæ occurrant, sumitur Conclusio Officii currentis, aut secus Officii ipsa die et primo quidem loco inter cetera propriam Conclusionem habentia commemorandi, aut demum de occurrenti Octava communi.}

\rubric{Expleto Hymno, conveniens dicitur Antiphona usque ad Asteriscum *, prouti occurrens Officium requirit.}

\variantheading{In Officio dominicali}

\rubric{In omnibus Dominicis, extra Octavas Ssmi Corporis Christi et Sacrat. Cordis Jesu, Antiphona ut in Psalterio.}

\variantheading{In Officio feriali}

\rubric{In omnibus Feriis, Antiphona de Feria currenti ut in Psalterio.}

\rubric{Post enuntiatam Antiphonam dicuntur tres Psalmi Officii diei congruentes.}

\rubric{¶ Quando autem adhibetur posterius Laudum schema, Psalmus, in priori Laudum ferialium schemate omissus, ad Primam Feriæ currentis post alios Psalmos resumitur.}

\rubric{Repetita post ultimum Psalmum integre Antiphona, conveniens dicitur Capitulum; nimirum:}

\rubric{In omnibus Dominicis, in Officio cujuslibet Festi vel Octavæ ac sanctæ Mariæ in Sabbato:}

\capitulumlabel{Capitulum}{1 Tim. 1, 17}

\lettrine[lines=2]{R}{EGI} sæculórum immortáli et invisíbili, soli Deo honor et glória in sæcula sæculórum. Amen. \R\ Deo grátias.

\rubric{¶ Sic respondetur in fine omnium Capitulorum.}

\rubric{In omnibus autem Feriis et in Vigiliis communibus:}

\capitulumlabel{Capitulum}{Zach. 8, 19}

\lettrine[lines=2]{P}{ACEM} et veritátem dilígite, ait Dóminus omnípotens.

\rubric{Dicto Capitulo, subjungitur Responsorium breve, ut sequitur. Versus tamen Qui sedes ad dexteram Patris semper omittitur, quando specialis Versus adhibendus præscribitur, Officii nempe currentis, aut secus Officii ad Laudes et primo quidem loco inter cetera specialem Versum habentia commemorati, aut demum de occurrenti Octava communi proprius.}

\R\ br. Christe, Fili Dei vivi, * Miserére nobis. \rubric{Et repetitur:} Christe, Fili Dei vivi, miserére nobis. \V\ Qui sedes ad déxteram Patris. \R\ Miserére nobis. \V\ Glória Patri, et Fílio, et Spirítui Sancto. \R\ Christe, Fili Dei vivi, miserére nobis. \V\ Exsúrge, Christe, ádjuva nos. \R\ Et líbera nos propter nomen tuum.

\rubric{¶ Ita dicuntur semper Responsoria brevia, etiam ad Tertiam, Sextam, Nonam et Completorium.}

\rubric{Expleto Responsorio brevi, si Preces non fuerint recitandæ, statim dicitur Versus Dóminus vobíscum cum subsequenti Oratione et reliquis omnibus usque ad finem Horæ ut infra.}

\rubric{In quolibet Officio semiduplici, in Officio sanctæ Mariæ in Sabbato, in Festis simplicibus, atque in Feriis per Annum communibus, post Responsorium breve dicuntur sequentes Preces dominicales; quæ tamen omittuntur omnibus diebus, in quibus ad Laudes habita sit Commemoratio Duplicis, vel factum fuerit Officium aut Commemoratio cujusvis Octavæ:}

KYRIE, eléison. Christe, eléison. Kýrie, eléison.

\rubric{Pater noster secreto usque ad} \V\ Et ne nos indúcas in tentatiónem. \R\ Sed líbera nos a malo.

\rubric{Credo in Deum secreto usque ad} \V\ Carnis resurrectiónem. \R\ Vitam ætérnam. Amen.

\V\ Et ego ad te, Dómine, clamávi. \R\ Et mane orátio mea prævéniet te. \V\ Repleátur os meum laude. \R\ Ut cantem glóriam tuam, tota die magnitúdinem tuam. \V\ Dómine, avérte fáciem tuam a peccátis meis. \R\ Et omnes iniquitátes meas dele. \V\ Cor mundum crea in me, Deus. \R\ Et spíritum rectum ínnova in viscéribus meis. \V\ Ne projícias me a fácie tua. \R\ Et spíritum sanctum tuum ne áuferas a me. \V\ Redde mihi lætítiam salutáris tui. \R\ Et spíritu principáli confírma me. \V\ Adjutórium nostrum in nómine Dómini. \R\ Qui fecit cælum et terram.

\rubric{Deinde Hebdomadarius facit Confessionem cum Absolutione, eamque Chorus repetit, ut infra post feriales Preces.}

\rubric{In omnibus vero Vigiliis, si ad Laudes recitatæ fuerint preces, post Responsorium breve dicuntur flexis genibus sequentes Preces feriales:}

KYRIE, eléison. Christe, eléison. Kýrie, eléison.

\rubric{Pater noster secreto usque ad} \V\ Et ne nos indúcas in tentatiónem. \R\ Sed líbera nos a malo.

\rubric{Credo in Deum secreto usque ad} \V\ Carnis resurrectiónem. \R\ Vitam ætérnam. Amen.

\V\ Et ego ad te, Dómine, clamávi. \R\ Et mane orátio mea prævéniet te. \V\ Repleátur os meum laude. \R\ Ut cantem glóriam tuam, tota die magnitúdinem tuam. \V\ Dómine, avérte fáciem tuam a peccátis meis. \R\ Et omnes iniquitátes meas dele. \V\ Cor mundum crea in me, Deus. \R\ Et spíritum rectum ínnova in viscéribus meis. \V\ Ne projícias me a fácie tua. \R\ Et spíritum sanctum tuum ne áuferas a me. \V\ Redde mihi lætítiam salutáris tui. \R\ Et spíritu principáli confírma me. \V\ Eripe me, Dómine, ab hómine malo. \R\ A viro iníquo éripe me. \V\ Eripe me de inimícis meis, Deus meus. \R\ Et ab insurgéntibus in me líbera me. \V\ Eripe me de operántibus iniquitátem. \R\ Et de viris sánguinum salva me. \V\ Sic psalmum dicam nómini tuo in sæculum sæculi. \R\ Ut reddam vota mea de die in diem. \V\ Exáudi nos, Deus, salutáris noster. \R\ Spes ómnium fínium terræ, et in mari longe. \V\ Deus, in adjutórium meum inténde. \R\ Dómine, ad adjuvándum me festína. \V\ Sanctus Deus, Sanctus fortis, Sanctus immortális. \R\ Miserére nobis. \V\ Bénedic, ánima mea, Dómino. \R\ Et ómnia, quæ intra me sunt, nómini sancto ejus. \V\ Bénedic, ánima mea, Dómino. \R\ Et noli oblivísci omnes retributiónes ejus. \V\ Qui propitiátur ómnibus iniquitátibus tuis. \R\ Qui sanat omnes infirmitátes tuas. \V\ Qui rédimit de intéritu vitam tuam. \R\ Qui corónat te in misericórdia et miseratiónibus. \V\ Qui replet in bonis desidérium tuum. \R\ Renovábitur ut áquilæ juvéntus tua. \V\ Adjutórium nostrum in nómine Dómini. \R\ Qui fecit cælum et terram.

\rubric{Deinde Hebdomadarius facit Confessionem:}

\lettrine[lines=2]{C}{ONFITEOR} Deo omnipoténti, beátæ Maríæ semper Vírgini, beáto Michaéli Archángelo, beáto Joánni Baptístæ, sanctis Apóstolis Petro et Paulo, ómnibus Sanctis, et vobis, fratres, quia peccávi nimis cogitatióne, verbo et ópere: mea culpa, mea culpa, mea máxima culpa. Ideo precor beátam Maríam semper Vírginem, beátum Michaélem Archángelum, beátum Joánnem Baptístam, sanctos Apóstolos Petrum et Paulum, omnes Sanctos, et vos, fratres, oráre pro me ad Dóminum, Deum nostrum.

\rubric{Chorus respondet:}

\lettrine[lines=2]{M}{ISEREATUR} tui omnípotens Deus, et, dimíssis peccátis tuis, perdúcat te ad vitam ætérnam. \R\ Amen.

\rubric{Deinde repetit Confessionem, et ubi dicitur: vobis, fratres, et vos, fratres, dicat: tibi, pater, et te, pater.}

\rubric{Facta Confessione a Choro, Hebdomadarius dicit:}

\lettrine[lines=2]{M}{ISEREATUR} vestri omnípotens Deus, et, dimíssis peccátis vestris, perdúcat vos ad vitam ætérnam. \R\ Amen.

\lettrine[lines=2]{I}{NDULGENTIAM}, absolutiónem, et remissiónem peccatórum nostrórum tríbuat nobis omnípotens et miséricors Dóminus. \R\ Amen.

\rubric{Extra Chorum, si unus vel duo tantum recitent Officium et in Choro Monialium, semel tantum ac simul ab omnibus ita fit Confessio:}

\lettrine[lines=2]{C}{ONFITEOR} Deo omnipoténti, beátæ Maríæ semper Vírgini, beáto Michaéli Archángelo, beáto Joánni Baptístæ, sanctis Apóstolis Petro et Paulo, et ómnibus Sanctis, quia peccávi nimis cogitatióne, verbo et ópere: mea culpa, mea culpa, mea máxima culpa. Ideo precor beátam Maríam semper Vírginem, beátum Michaélem Archángelum, beátum Joánnem Baptístam, sanctos Apóstolos Petrum et Paulum, et omnes Sanctos, oráre pro me ad Dóminum, Deum nostrum.

\rubric{Deinde dicitur:}

\lettrine[lines=2]{M}{ISEREATUR} nostri omnípotens Deus, et, dimíssis peccátis nostris, perdúcat nos ad vitam ætérnam. Amen.

\lettrine[lines=2]{I}{NDULGENTIAM}, absolutiónem, et remissiónem peccatórum nostrórum tríbuat nobis omnípotens et miséricors Dóminus. Amen.

\rubric{Facta Absolutione, subjungitur:}

\V\ Dignáre, Dómine, die isto. \R\ Sine peccáto nos custodíre. \V\ Miserére nostri, Dómine. \R\ Miserére nostri. \V\ Fiat misericórdia tua, Dómine, super nos. \R\ Quemádmodum sperávimus in te. \V\ Dómine, exáudi oratiónem meam. \R\ Et clamor meus ad te véniat.

\rubric{Absolutis Precibus, vel, si Preces locum non habeant, post Responsorium breve statim dicitur:}

\V\ Dóminus vobíscum. \R\ Et cum spíritu tuo.

Orémus. \oratiolabel{Oratio}

\lettrine[lines=2]{D}{OMINE}, Deus omnípotens, qui ad princípium hujus diéi nos perveníre fecísti: tua nos hódie salva virtúte; ut in hac die ad nullum declinémus peccátum, sed semper ad tuam justítiam faciéndam nostra procédant elóquia, dirigántur cogitatiónes et ópera. Per Dóminum. \R\ Amen.

\V\ Dóminus vobíscum. \R\ Et cum spíritu tuo. \V\ Benedicámus Dómino. \R\ Deo grátias.

\rubric{Deinde in Choro legitur Martyrologium, quod laudabiliter fit etiam extra Chorum. Postea Hebdomadarius dicit:}

\V\ Pretiósa in conspéctu Dómini. \R\ Mors Sanctórum ejus.

\rubric{Et absolute, sine} Orémus:

\oratiolabel{Oratio}

\lettrine[lines=2]{S}{ANCTA} María, et omnes Sancti intercédant pro nobis ad Dóminum, ut nos mereámur ab eo adjuvári et salvári, qui vivit et regnat in sæcula sæculórum. \R\ Amen.

\V\ Deus, in adjutórium meum inténde. \R\ Dómine, ad adjuvándum me festína. \rubric{Et dicitur ter; ultimo additur:} Glória Patri, et Fílio, et Spirítui Sancto. Sicut erat in princípio, et nunc, et semper, et in sæcula sæculórum. Amen.

Kýrie, eléison. Christe, eléison. Kýrie, eléison.

Pater noster \rubric{secreto usque ad} \V\ Et ne nos indúcas in tentatiónem. \R\ Sed líbera nos a malo.

\V\ Réspice in servos tuos, Dómine, et in ópera tua, et dírige fílios eórum. \R\ Et sit splendor Dómini, Dei nostri, super nos, et ópera mánuum nostrárum dírige super nos, et opus mánuum nostrárum dírige. \V\ Glória Patri, et Fílio, et Spirítui Sancto. \R\ Sicut erat in princípio, et nunc, et semper, et in sæcula sæculórum. Amen.

Orémus. \oratiolabel{Oratio}

\lettrine[lines=2]{D}{IRIGERE} et sanctificáre, régere et gubernáre dignáre, Dómine Deus, Rex cæli et terræ, hódie corda et córpora nostra, sensus, sermónes et actus nostros in lege tua, et in opéribus mandatórum tuórum: ut hic, et in ætérnum, te auxiliánte, salvi et líberi esse mereámur, Salvátor mundi: Qui vivis et regnas in sæcula sæculórum. \R\ Amen.

\rubric{Deinde dicitur Lectio brevis, ad absolutionem Capituli, prouti Officium occurrens requirit.}

\V\ Jube, domne, benedícere. \benedlabel{Bened.} Dies et actus nostros in sua pace dispónat Dóminus omnípotens. \R\ Amen.

\variantheading{In Officio dominicali et feriali}

\rubric{A Feria II usque ad Feriam IV post Festum sanctissimæ Trinitatis, et a Sabbato post Octavam Sacratissimi Cordis Jesu et deinceps.}

\lectioheading{Lectio brevis}{2 Thess. 3, 5}

\lettrine[lines=2]{D}{OMINUS} autem dírigat corda et córpora nostra in caritáte Dei, et patiéntia Christi. Tu autem, Dómine, miserére nobis. \R\ Deo grátias.

\rubric{¶ Ita semper concluduntur Lectiones breves ad Primam, et ita in earum fine respondetur.}

\rubric{In Officio autem cujuslibet Festi vel Octavæ, atque S. Mariæ in Sabbato; itemque in Dominicis infra Octavas Ssmi Corporis Christi ac Sacrat. Cordis Jesu, si de eis fiat Officium:}

\rubric{Pro Lectione brevi sumitur Capitulum, quod ad Nonam ponitur ipsius Officii currentis, et in fine dicitur: Tu autem, Dómine, ut supra.}

\rubric{Post Lectionem brevem dicitur:}

\V\ Adjutórium nostrum in nómine Dómini. \R\ Qui fecit cælum et terram.

\V\ Benedícite. \R\ Deus.

\benedlabel{Bened.} Dóminus nos benedícat, et ab omni malo deféndat, et ad vitam perdúcat ætérnam. Et fidélium ánimæ per misericórdiam Dei requiéscant in pace. \R\ Amen.

\rubric{Deinde dicitur tantum Pater noster secreto, nisi sequatur alia Hora.}

\hourheading{Ad Tertiam}

Pater noster et Ave.

\V\ Deus, in adjutórium.

\hymnusheading{Hymnus}

\lettrine[lines=2]{N}{UNC}, Sancte, nobis, Spíritus,\\
Unum Patri cum Fílio,\\
Dignáre promptus íngeri\\
Nostro refúsus péctori.

Os, lingua, mens, sensus, vigor,\\
Confessiónem pérsonent,\\
Flamméscat igne cáritas,\\
Accéndat ardor próximos.

Præsta, Pater piíssime,\\
Patríque compar Unice,\\
Cum Spíritu Paráclito\\
Regnans per omne sæculum. Amen.

\rubric{Expleto Hymno, conveniens dicitur Antiphona usque ad Asteriscum *, et tres Psalmi, prouti Officium occurrens requirit; atque Psalmis absolutis, integra repetitur Antiphona.}

\variantheading{In Officio dominicali}

\rubric{In omnibus Dominicis, extra Octavas Ssmi Corporis Christi ac Sacrat. Cordis Jesu, Aña ut in Psalterio, et dicitur sequens}

\capitulumlabel{Capitulum}{1 Joann. 4, 16}

\lettrine[lines=2]{D}{EUS} cáritas est: et qui manet in caritáte, in Deo manet, et Deus in eo.

\R\ br. Inclína cor meum, Deus, * In testimónia tua. Inclína. \V\ Avérte óculos meos, ne vídeant vanitátem: in via tua vivífica me. In testimónia tua. Glória Patri. Inclína.

\V\ Ego dixi: Dómine, miserére mei. \R\ Sana ánimam meam, quia peccávi tibi.

\variantheading{In Officio feriali}

\rubric{In omnibus Feriis, Antiphona de Feria currenti, ut in Psalterio; et dicitur quotidie sequens}

\capitulumlabel{Capitulum}{Jer. 17, 14}

\lettrine[lines=2]{S}{ANA} me, Dómine, et sanábor: salvum me fac, et salvus ero: quóniam laus mea tu es.

\R\ br. Sana ánimam meam, * Quia peccávi tibi. Sana. \V\ Ego dixi: Dómine, miserére mei. Quia. Glória Patri. Sana.

\V\ Adjútor meus esto, ne derelínquas me. \R\ Neque despícias me, Deus, salutáris meus.

\rubric{Expleto Responsorio brevi, si Preces non fuerint recitandæ, statim dicitur Versus Dóminus vobíscum cum Oratione et reliquis omnibus usque ad finem Horæ, ut infra.}

\rubric{Quando vero ad Laudes recitatæ sunt Preces, post Responsorium breve dicuntur sequentes flexis genibus:}

KYRIE, eléison. Christe, eléison. Kýrie, eléison.

Pater noster \rubric{secreto usque ad} \V\ Et ne nos indúcas in tentatiónem. \R\ Sed líbera nos a malo. \V\ Dómine, Deus virtútum, convérte nos. \R\ Et osténde fáciem tuam, et salvi érimus. \V\ Exsúrge, Christe, ádjuva nos. \R\ Et líbera nos propter nomen tuum. \V\ Dómine, exáudi oratiónem meam. \R\ Et clamor meus ad te véniat.

\rubric{Absolutis Precibus, vel, si Preces locum non habeant, post Responsorium breve statim dicitur:}

\V\ Dóminus vobíscum. \R\ Et cum spíritu tuo.

Orémus.

\rubric{Et dicitur Oratio conveniens.}

\rubric{Post Orationem additur:}

\V\ Dóminus vobíscum. \R\ Et cum spíritu tuo.

\V\ Benedicámus Dómino. \R\ Deo grátias.

\V\ Fidélium ánimæ per misericórdiam Dei requiéscant in pace. \R\ Amen.

\rubric{Deinde dicitur tantum Pater noster secreto, nisi sequatur alia Hora.}

\hourheading{Ad Sextam}

Pater noster et Ave.

\V\ Deus, in adjutórium.

\hymnusheading{Hymnus}

\lettrine[lines=2]{R}{ECTOR} potens, verax Deus,\\
Qui témperas rerum vices,\\
Splendóre mane illúminas,\\
Et ígnibus merídiem:

Exstíngue flammas lítium,\\
Aufer calórem nóxium,\\
Confer salútem córporum,\\
Verámque pacem córdium.

Præsta, Pater piíssime,\\
Patríque compar Unice,\\
Cum Spíritu Paráclito\\
Regnans per omne sæculum. Amen.

\rubric{Expleto Hymno, conveniens dicitur Antiphona usque ad Asteriscum *, et tres Psalmi prouti Officium occurrens requirit; atque, Psalmis absolutis, integra repetitur Antiphona.}

\variantheading{In Officio dominicali}

\rubric{In omnibus Dominicis, extra Octavas Ssmi Corporis Christi ac Sacrat. Cordis Jesu, Aña ut in Psalterio, et dicitur sequens}

\capitulumlabel{Capitulum}{Gal. 6, 2}

\lettrine[lines=2]{A}{LTER} altérius ónera portáte, et sic adimplébitis legem Christi.

\R\ br. In ætérnum, Dómine, * Pérmanet verbum tuum. In ætérnum. \V\ In sæculum sæculi véritas tua. Pérmanet. Glória Patri. In ætérnum.

\V\ Dóminus regit me, et nihil mihi déerit. \R\ In loco páscuæ ibi me collocávit.

\variantheading{In Officio feriali}

\rubric{In omnibus Feriis, Antiphona de Feria currenti, ut in Psalterio; et dicitur quotidie sequens}

\capitulumlabel{Capitulum}{Rom. 13, 8}

\lettrine[lines=2]{N}{EMINI} quidquam debeátis, nisi ut ínvicem diligátis: qui enim díligit próximum, legem implévit.

\R\ br. Benedícam Dóminum * In omni témpore. Benedícam. \V\ Semper laus ejus in ore meo. In. Glória Patri. Benedícam.

\V\ Dóminus regit me, et nihil mihi déerit. \R\ In loco páscuæ ibi me collocávit.

\rubric{Expleto Responsorio brevi, si Preces non fuerint recitandæ, statim dicitur Versus Dóminus vobíscum cum Oratione et reliquis omnibus usque ad finem Horæ, ut infra.}

\rubric{Quando vero ad Laudes recitatæ sunt Preces, post Responsorium breve dicuntur sequentes flexis genibus:}

KYRIE, eléison. Christe, eléison. Kýrie, eléison.

Pater noster \rubric{secreto usque ad} \V\ Et ne nos indúcas in tentatiónem. \R\ Sed líbera nos a malo. \V\ Dómine, Deus virtútum, convérte nos. \R\ Et osténde fáciem tuam, et salvi érimus. \V\ Exsúrge, Christe, ádjuva nos. \R\ Et líbera nos propter nomen tuum. \V\ Dómine, exáudi oratiónem meam. \R\ Et clamor meus ad te véniat.

\rubric{Absolutis Precibus, vel, si Preces locum non habeant, post Responsorium breve statim dicitur:}

\V\ Dóminus vobíscum. \R\ Et cum spíritu tuo.

Orémus.

\rubric{Et dicitur Oratio conveniens.}

\rubric{Post Orationem additur:}

\V\ Dóminus vobíscum. \R\ Et cum spíritu tuo.

\V\ Benedicámus Dómino. \R\ Deo grátias.

\V\ Fidélium ánimæ per misericórdiam Dei requiéscant in pace. \R\ Amen.

\rubric{Deinde dicitur tantum Pater noster secreto, nisi sequatur alia Hora.}

\hourheading{Ad Nonam}

Pater noster et Ave.

\V\ Deus, in adjutórium.

\hymnusheading{Hymnus}

\lettrine[lines=2]{R}{ERUM}, Deus, tenax vigor,\\
Immótus in te pérmanens,\\
Lucis diúrnæ témpora\\
Successíbus detérminans:

Largíre lumen véspere,\\
Quo vita nusquam décidat,\\
Sed præmium mortis sacræ\\
Perénnis instet glória.

Præsta, Pater piíssime,\\
Patríque compar Unice,\\
Cum Spíritu Paráclito\\
Regnans per omne sæculum. Amen.

\rubric{Expleto Hymno, conveniens dicitur Antiphona usque ad Asteriscum *, et tres Psalmi, prouti Officium occurrens requirit; atque, Psalmis absolutis, integra repetitur Antiphona.}

\variantheading{In Officio dominicali}

\rubric{In omnibus Dominicis, extra Octavas Ssmi Corporis Christi ac Sacrat. Cordis Jesu, Aña ut in Psalterio, et dicitur sequens}

\capitulumlabel{Capitulum}{1 Cor. 6, 20}

\lettrine[lines=2]{E}{MPTI} enim estis prétio magno. Glorificáte et portáte Deum in córpore vestro.

\R\ br. Clamávi in toto corde meo: * Exáudi me, Dómine. Clamávi. \V\ Justificatiónes tuas requíram. Exáudi. Glória Patri. Clamávi.

\V\ Ab occúltis meis munda me, Dómine. \R\ Et ab aliénis parce servo tuo.

\variantheading{In Officio feriali}

\rubric{In omnibus Feriis, Antiphona de Feria currenti, ut in Psalterio; et dicitur quotidie sequens}

\capitulumlabel{Capitulum}{1 Petr. 1, 17-19}

\lettrine[lines=2]{I}{N} timóre incolátus vestri témpore conversámini: sciéntes quod non corruptibílibus auro vel argénto redémpti estis, sed pretióso sánguine quasi Agni immaculáti Christi.

\R\ br. Rédime me, Dómine, * Et miserére mei. Rédime. \V\ Pes enim meus stetit in via recta. Et. Glória Patri. Rédime.

\V\ Ab occúltis meis munda me, Dómine. \R\ Et ab aliénis parce servo tuo.

\rubric{Expleto Responsorio brevi, si Preces non fuerint recitandæ, statim dicitur Versus Dóminus vobíscum cum Oratione et reliquis omnibus usque ad finem Horæ, ut infra.}

\rubric{Quando vero ad Laudes recitatæ sunt Preces, post Responsorium breve dicuntur sequentes flexis genibus:}

KYRIE, eléison. Christe, eléison. Kýrie, eléison.

Pater noster \rubric{secreto usque ad} \V\ Et ne nos indúcas in tentatiónem. \R\ Sed líbera nos a malo. \V\ Dómine, Deus virtútum, convérte nos. \R\ Et osténde fáciem tuam, et salvi érimus. \V\ Exsúrge, Christe, ádjuva nos. \R\ Et líbera nos propter nomen tuum. \V\ Dómine, exáudi oratiónem meam. \R\ Et clamor meus ad te véniat.

\rubric{Absolutis Precibus, vel, si Preces locum non habeant, post Responsorium breve statim dicitur:}

\V\ Dóminus vobíscum. \R\ Et cum spíritu tuo.

Orémus.

\rubric{Et dicitur Oratio conveniens.}

\rubric{Post Orationem additur:}

\V\ Dóminus vobíscum. \R\ Et cum spíritu tuo.

\V\ Benedicámus Dómino. \R\ Deo grátias.

\V\ Fidélium ánimæ per misericórdiam Dei requiéscant in pace. \R\ Amen.

\rubric{Deinde dicitur tantum Pater noster secreto, nisi sequatur alia Hora.}

% ===== END OF BATCH 2 -- Covers aestiva.md lines 465-874, completing
% "Ad Primam" (hymn "Jam lucis orto sidere," the full Confiteor/
% Misereatur/Indulgentiam apparatus, the Martyrologium cue, the
% "Dirigere et sanctificare" Oratio, and the Lectio brevis, here
% given in its single simplified Pars Aestiva form -- "In Officio
% dominicali et feriali" -- rather than the three-way Lenten/
% Paschal-season branching found in verna.tex's Prime, since this
% volume's Lectio brevis never varies by season), "Ad Tertiam," "Ad
% Sextam," and "Ad Nonam," each with its proper hymn and its
% Sunday/weekday Capitulum-and-brief-responsory pair. Two spellings
% of the same word recur within this single batch exactly as the
% source alternates them: "sæculum"/"sæculi" at Sexta's responsory
% (line 758) versus the different spelling pattern seen in the Te
% Deum of Batch 1 -- each transcribed as its own occurrence gives it,
% per the standing rule against silent normalization. The document
% continues with "Ad Vesperas" (aestiva.md line 875) =====
